\documentclass{article}

\usepackage[a4paper,margin=1in]{geometry}

\usepackage{amsmath, amsthm, amssymb, amsfonts}
\usepackage{mathtools}
\usepackage{enumitem}
\usepackage{microtype}
\usepackage{xcolor}
\usepackage{fancyhdr}
\usepackage{graphicx}
\usepackage{hyperref}

% -------------------------------
% Fonts and Encoding
% -------------------------------
\usepackage{lmodern}
\usepackage[T1]{fontenc}
\usepackage[utf8]{inputenc}

% -------------------------------
% Theorem Environments
% -------------------------------
\theoremstyle{definition}
\newtheorem{definition}{Definition}[section]
\newtheorem{example}[definition]{Example}
\newtheorem{remark}[definition]{Remark}

\theoremstyle{plain}
\newtheorem{theorem}[definition]{Theorem}
\newtheorem{lemma}[definition]{Lemma}
\newtheorem{corollary}[definition]{Corollary}
\newtheorem{proposition}[definition]{Proposition}

% -------------------------------
% Header and Footer
% -------------------------------
\pagestyle{fancy}
\fancyhf{}
\rhead{\thepage}
\lhead{\textit{Mathematical Notes}}
\cfoot{}

% -------------------------------
% Custom Commands
% -------------------------------
\newcommand{\R}{\mathbb{R}}
\newcommand{\Z}{\mathbb{Z}}
\newcommand{\Q}{\mathbb{Q}}
\newcommand{\N}{\mathbb{N}}
\newcommand{\C}{\mathbb{C}}

\DeclareMathOperator{\Ker}{Ker}
\DeclareMathOperator{\Img}{Im}

% -------------------------------
% Examples & exercises
% -------------------------------
\newenvironment{solution}{\par\noindent\textbf{Solution.}\ }{\par}
\newenvironment{exercise}{\par\noindent\textbf{Exercise.}\ }{\par}

\begin{document}
	
	\title{My Mathematics Notes}
	\author{Ramon}
	\date{\today}
	
	\maketitle
	The ring $(S,+,\cdot)$ with 
	\[
	1 \neq 0,
	\]
	\[
	ab=ba \qquad \forall a,b \in S,
	\]
	and
	\[
	\forall a \in S,\quad a \neq 0 \implies \exists a^{-1} \in S \text{ such that } aa^{-1}=a^{-1}a=1.
	\]
	is called \textbf{field}
	
	Let $S$ be a field. The field $S$ already governs how scalars interact with scalars through its own addition and multiplication. Separately, let $V$ be a set equipped with an addition
	\[
	+ : V\times V \to V
	\]
	such that $(V,+)$ is an abelian group. Thus $V$ provides the abstract vectors together with their internal additive structure.
	
	To connect the scalar structure of $S$ with the vector structure of $V$, one imposes an external operation
	\[
	\cdot : S\times V \to V,
	\]
	usually written $(a,v)\mapsto av$, which tells how a scalar acts on a vector.
	
	A \textbf{vector space} over $S$ is therefore a set $V$ equipped with maps
	\[
	+ : V\times V \to V,
	\qquad
	\cdot : S\times V \to V,
	\]
	such that $(V,+)$ is an abelian group and, for all $a,b\in S$ and $u,v\in V$,
	\[
	a(u+v)=au+av,\qquad (a+b)v=av+bv,\qquad (ab)v=a(bv),\qquad 1v=v.
	\]
	A \textbf{linear combination} of vectors \(v_1,\dots,v_n\in V\) is an expression of the form
	\[
	a_1v_1+\cdots+a_nv_n,
	\qquad a_1,\dots,a_n\in S.
	\]
	It is formed by first scaling each vector \(v_i\) by a coefficient \(a_i\), and then adding the resulting vectors in \(V\).
	
	Unlike the ordinary sum
	\[
	v_1+\cdots+v_n,
	\]
	which corresponds to taking every coefficient equal to \(1\), a linear combination allows one not only to combine vectors, but also to control the contribution of each vector through its coefficient. In this way, scalar multiplication together with vector addition gives rise to the basic structure of a linear combination.
	Thus a vector space is the structure in which scalars from $S$ act coherently on vectors in $V$, making possible expressions of the form
	\[
	a_1v_1+\cdots+a_nv_n.
	\]
	Hence vector spaces are about \textbf{linear combinations of vectors with scalar coefficients}, not about combining scalars alone. In the expression
	\[
	a_1v_1+\cdots+a_nv_n,
	\]
	each $v_i\in V$ provides a vector, and each scalar $a_i\in S$ provides its coefficient. The role of $a_i$ is to determine the scaled contribution of $v_i$ to the total combination. Thus $v_i$ selects the vector, while $a_i$ determines its weight in the sum.
	
	Once linear combinations
	\[
	a_1v_1+\cdots+a_nv_n
	\]
	are available, one may choose a family of vectors
	\[
	v_1,\dots,v_n\in V
	\]
	and ask whether this family contains redundancy. Redundancy means that one of the vectors adds no genuinely new linear contribution because it can already be obtained from the others by a linear combination.
	
	This leads to the notion of linear dependence. The family $v_1,\dots,v_n$ is said to be \textbf{linearly dependent} if there exist scalars $a_1,\dots,a_n\in S$, not all zero, such that
	\[
	a_1v_1+\cdots+a_nv_n=0.
	\]
	Such a relation is called a nontrivial linear relation among the vectors. It shows that the family contains redundancy.
	
	Equivalently, if one of the coefficients, say $a_k$, is nonzero, then one may solve for $v_k$ as a linear combination of the remaining vectors. Hence one of the vectors is already determined by the others and does not contribute a new independent direction.
	
	By contrast, the family $v_1,\dots,v_n$ is \textbf{linearly independent} if
	\[
	a_1v_1+\cdots+a_nv_n=0
	\quad\Longrightarrow\quad
	a_1=\cdots=a_n=0.
	\]
	Thus linear independence means that the only linear combination of the chosen vectors which gives the zero vector is the trivial one. In this case, the family contains no linear redundancy.
	
	Once the vector space structure has been established and linear combinations are available, one may choose a family of vectors in \(V\). To organize this family, one introduces an index set \(I\) and writes
	\[
	\{v_i\}_{i\in I}\subseteq V.
	\]
	This indexed family simply labels the chosen vectors of \(V\); to refer to them cleanly 
	
	Given a family of vectors
	\[
	\{v_i\}_{i\in I}\subseteq V,
	\]
	one may ask whether every vector in the family is genuinely needed. The family is said to contain \textbf{redundancy} if one of its vectors can be expressed as a linear combination of the others. In that case, that vector adds no new linear contribution, since the same generated vectors could be obtained without it. This idea leads to the formal notion of linear dependence.
	
	A family of vectors
	\[
	\{v_i\}_{i\in I}\subseteq V
	\]
	is said to be \textbf{linearly dependent} if there exist finitely many indices
	\[
	i_1,\dots,i_n\in I
	\]
	and coefficients
	\[
	a_1,\dots,a_n\in S,
	\]
	not all zero, such that
	\[
	a_1v_{i_1}+\cdots+a_nv_{i_n}=0.
	\]
	Thus linear dependence means that the family satisfies a nontrivial linear relation. Equivalently, at least one vector of the family can be expressed as a linear combination of the others, and hence the family contains redundancy.
	
	As an illustration, consider the vector space \(\mathbb{R}^2\).
	
	First, let
	\[
	u_1=(1,0),\qquad u_2=(0,1),\qquad u_3=(1,1).
	\]
	This family contains redundancy, since
	\[
	u_3=u_1+u_2.
	\]
	Equivalently,
	\[
	u_1+u_2-u_3=0,
	\]
	and the coefficients \(1,1,-1\) are not all zero. Hence the family
	\[
	\{u_1,u_2,u_3\}
	\]
	is linearly dependent. The vector \(u_3\) contributes no new linear information, since it can already be obtained from the others.
	
	Now consider instead the family
	\[
	e_1=(1,0),\qquad e_2=(0,1).
	\]
	Suppose
	\[
	a_1e_1+a_2e_2=0.
	\]
	Then
	\[
	a_1(1,0)+a_2(0,1)=(0,0),
	\]
	so
	\[
	(a_1,a_2)=(0,0),
	\]
	which implies
	\[
	a_1=0,\qquad a_2=0.
	\]
	Therefore the family
	\[
	\{e_1,e_2\}
	\]
	is linearly independent.
	
	Thus the first family is linearly dependent because it contains redundancy, whereas the second is linearly independent because no vector in it can be obtained as a linear combination of the others.
	
	Given a family of vectors
	\[
	\{v_i\}_{i\in I}\subseteq V,
	\]
	the \textbf{span} of this family is the set of all finite linear combinations of its vectors:
	\[
	\operatorname{Span}\{v_i\}_{i\in I}
	=
	\left\{
	a_1v_{i_1}+\cdots+a_nv_{i_n}
	\;\middle|\;
	n\geq 1,\;
	i_1,\dots,i_n\in I,\;
	a_1,\dots,a_n\in S
	\right\}.
	\]
	Thus the span of a family is the set of all vectors that can be generated from it by scalar multiplication and vector addition. If
	\[
	\operatorname{Span}\{v_i\}_{i\in I}=V,
	\]
	then the family is said to \textbf{span} the vector space \(V\). The notion of span does not require linear independence. Any family of vectors has a span, namely the set of all finite linear combinations of its elements. Linear independence is an additional condition which excludes redundancy. A basis is therefore a family of vectors which spans the space and is also linearly independent.
	
	Let \(V\) be a vector space over \(S\), and let
	\[
	\{v_i\}_{i\in I}\subseteq V
	\]
	be a family of vectors. This family is called a \textbf{basis} of \(V\) if it satisfies the following two conditions:
	
	\[
	\operatorname{Span}\{v_i\}_{i\in I}=V,
	\]
	that is, every vector in \(V\) can be expressed as a finite linear combination of elements of the family; and
	
	\[
	\{v_i\}_{i\in I}\text{ is linearly independent},
	\]
	that is, the only linear combination of its elements equal to \(0\) is the trivial one.
	
	Thus a basis is a family of vectors that generates the whole space without redundancy. Equivalently, a basis is a linearly independent spanning family of \(V\).
	
	Consider the vector space
	\[
	V=\mathbb{R}^2
	\]
	over the field
	\[
	S=\mathbb{R}.
	\]
	Its elements are vectors
	\[
	v=(x,y)\in V.
	\]
	
	Let us choose the family
	\[
	u_1=(1,0),\qquad u_2=(0,1),\qquad u_3=(1,1).
	\]
	This may be written as an indexed family
	\[
	\{u_i\}_{i=1}^3\subseteq V.
	\]
	
	A general linear combination of these vectors has the form
	\[
	a_1u_1+a_2u_2+a_3u_3,
	\qquad a_1,a_2,a_3\in\mathbb{R}.
	\]
	Substituting the vectors gives
	\[
	a_1(1,0)+a_2(0,1)+a_3(1,1)
	=
	(a_1+a_3,\; a_2+a_3).
	\]
	Thus the coefficients determine how much each vector contributes to the final combination.
	
	Now observe that
	\[
	u_3=u_1+u_2.
	\]
	Equivalently,
	\[
	u_1+u_2-u_3=0.
	\]
	Since the coefficients \(1,1,-1\) are not all zero, this is a nontrivial linear relation. Therefore the family
	\[
	\{u_1,u_2,u_3\}
	\]
	is linearly dependent. This shows that the family contains redundancy: the vector \(u_3\) adds no new linear information, because it can already be obtained from \(u_1\) and \(u_2\).
	
	Let us now compute its span. By definition,
	\[
	\operatorname{Span}\{u_1,u_2,u_3\}
	\]
	is the set of all vectors of the form
	\[
	a_1u_1+a_2u_2+a_3u_3.
	\]
	But we have already seen that every such combination is
	\[
	(a_1+a_3,\; a_2+a_3).
	\]
	Conversely, given any vector \((x,y)\in\mathbb{R}^2\), one may choose
	\[
	a_1=x,\qquad a_2=y,\qquad a_3=0,
	\]
	so that
	\[
	(x,y)=xu_1+yu_2.
	\]
	Hence every vector in \(\mathbb{R}^2\) belongs to the span, and therefore
	\[
	\operatorname{Span}\{u_1,u_2,u_3\}=\mathbb{R}^2.
	\]
	So this family spans the whole space, but it is not a basis, because it is linearly dependent.
	
	Now remove the redundant vector \(u_3\), and consider the family
	\[
	e_1=(1,0),\qquad e_2=(0,1).
	\]
	Then
	\[
	\{e_1,e_2\}\subseteq \mathbb{R}^2.
	\]
	
	Its span is again all of \(\mathbb{R}^2\), because every vector \((x,y)\in\mathbb{R}^2\) may be written as
	\[
	(x,y)=xe_1+ye_2.
	\]
	Thus
	\[
	\operatorname{Span}\{e_1,e_2\}=\mathbb{R}^2.
	\]
	
	We now check linear independence. Suppose
	\[
	a_1e_1+a_2e_2=0.
	\]
	Then
	\[
	a_1(1,0)+a_2(0,1)=(0,0),
	\]
	so
	\[
	(a_1,a_2)=(0,0).
	\]
	Hence
	\[
	a_1=0,\qquad a_2=0.
	\]
	Therefore the family \(\{e_1,e_2\}\) is linearly independent.
	
	We conclude that
	\[
	\{e_1,e_2\}
	\]
	is a basis of \(\mathbb{R}^2\), because it spans the whole space and contains no redundancy.
	
	This example shows the full chain of ideas:
	\[
	\text{family of vectors}
	\;\to\;
	\text{linear combinations}
	\;\to\;
	\text{redundancy}
	\;\to\;
	\text{linear dependence}
	\;\to\;
	\text{span}
	\;\to\;
	\text{basis}.
	\]
	
	Let \(V\) be a vector space over \(\mathbb{C}\), and let
	\[
	B=\{e_i\}_{i\in I}
	\]
	be a basis of \(V\). Since \(B\) spans \(V\), every vector \(v\in V\) can be written as a finite linear combination of basis vectors:
	\[
	v=a_1e_{i_1}+\cdots+a_ne_{i_n},
	\qquad a_1,\dots,a_n\in\mathbb{C}.
	\]
	Since \(B\) is also linearly independent, this representation is unique. Thus a basis provides a coordinate description of each vector: the vector itself remains an abstract element of \(V\), while the coefficients \(a_1,\dots,a_n\) describe that vector relative to the chosen basis.
	
	Let \(V\) be a vector space over \(\mathbb{C}\), and let
	\[
	B=\{e_i\}_{i\in I}
	\]
	be a basis of \(V\). Since \(B\) spans \(V\), every vector \(v\in V\) can be written as a finite linear combination of basis vectors:
	\[
	v=\sum_{k=1}^{n} a_k e_{i_k},
	\qquad a_k\in\mathbb{C},\ i_k\in I.
	\]
	Since \(B\) is also linearly independent, this representation is unique. Thus a basis provides a coordinate description of each vector: the vector itself remains an abstract element of \(V\), while the coefficients \(a_k\) describe that vector relative to the chosen basis.
	
	Let \(V\) and \(W\) be vector spaces over \(\mathbb{C}\). A map
	\[
	T:V\to W
	\]
	is called a \textbf{linear map} if, for all \(u,v\in V\) and all \(a\in\mathbb{C}\),
	\[
	T(u+v)=T(u)+T(v),
	\qquad
	T(av)=aT(v).
	\]
	Thus a linear map preserves the operations that define the vector-space structure: vector addition and multiplication by coefficients. Equivalently, for all \(u,v\in V\) and all \(a,b\in\mathbb{C}\),
	\[
	T(au+bv)=aT(u)+bT(v).
	\]
	Hence a linear map sends linear combinations in \(V\) to the corresponding linear combinations in \(W\).
	
	You start with two elements
	\[
	a,b\in G.
	\]
	
	There are then two possible routes.
	
	\medskip
	
	\textbf{Route 1: combine first in \(G\), then map.}
	
	First use the operation of \(G\):
	\[
	a\ast b\in G.
	\]
	
	Then apply \(\varphi\):
	\[
	\varphi(a\ast b)\in H.
	\]
	
	\medskip
	
	\textbf{Route 2: map first to \(H\), then combine in \(H\).}
	
	First apply \(\varphi\) separately:
	\[
	\varphi(a),\varphi(b)\in H.
	\]
	
	Then use the operation of \(H\):
	\[
	\varphi(a)\cdot\varphi(b)\in H.
	\]
	
	\medskip
	
	The homomorphism condition says that these two routes give the same result:
	\[
	\varphi(a\ast b)=\varphi(a)\cdot\varphi(b).
	\]
	
	You start with two rings
	\[
	(R,+,\cdot)\qquad\text{and}\qquad(S,\oplus,\otimes).
	\]
	
	A ring homomorphism is a function
	\[
	\varphi:R\to S
	\]
	such that it is compatible with both ring operations.
	
	\medskip
	
	\textbf{Addition.}
	
	You start with two elements
	\[
	a,b\in R.
	\]
	
	There are then two possible routes.
	
	\medskip
	
	\textbf{Route 1: add first in \(R\), then map.}
	
	First use the addition in \(R\):
	\[
	a+b\in R.
	\]
	
	Then apply \(\varphi\):
	\[
	\varphi(a+b)\in S.
	\]
	
	\medskip
	
	\textbf{Route 2: map first to \(S\), then add in \(S\).}
	
	First apply \(\varphi\) separately:
	\[
	\varphi(a),\varphi(b)\in S.
	\]
	
	Then use the addition in \(S\):
	\[
	\varphi(a)\oplus\varphi(b)\in S.
	\]
	
	\medskip
	
	The homomorphism condition for addition says that these two routes give the same result:
	\[
	\varphi(a+b)=\varphi(a)\oplus\varphi(b).
	\]
	
	\bigskip
	
	\textbf{Multiplication.}
	
	Again start with
	\[
	a,b\in R.
	\]
	
	There are two possible routes.
	
	\medskip
	
	\textbf{Route 1: multiply first in \(R\), then map.}
	
	First use the multiplication in \(R\):
	\[
	ab\in R.
	\]
	
	Then apply \(\varphi\):
	\[
	\varphi(ab)\in S.
	\]
	
	\medskip
	
	\textbf{Route 2: map first to \(S\), then multiply in \(S\).}
	
	First apply \(\varphi\) separately:
	\[
	\varphi(a),\varphi(b)\in S.
	\]
	
	Then use the multiplication in \(S\):
	\[
	\varphi(a)\otimes\varphi(b)\in S.
	\]
	
	\medskip
	
	The homomorphism condition for multiplication says that these two routes give the same result:
	\[
	\varphi(ab)=\varphi(a)\otimes\varphi(b).
	\]
	
	Thus, from the ring perspective, a homomorphism is a function that preserves both operations defining the ring structure: addition and multiplication.
	
	A ring homomorphism is a map
	\[
	\varphi:R\to S
	\]
	such that for all \(a,b\in R\),
	\[
	\varphi(a+b)=\varphi(a)+\varphi(b)
	\qquad\text{and}\qquad
	\varphi(ab)=\varphi(a)\varphi(b).
	\]
	
	Although a ring has two operations, a ring homomorphism has only one kernel:
	\[
	\ker(\varphi)=\{r\in R:\varphi(r)=0\}.
	\]
	
	This kernel is defined with respect to the additive structure, because every ring is, in particular, an abelian group under addition, and
	\[
	\varphi:(R,+)\to(S,+)
	\]
	is a group homomorphism.
	
	There is no second kernel coming from multiplication. Multiplication does not produce an independent kernel notion here, since the multiplicative structure is not, in general, a group structure on all of \(R\). Its role is different: it does not define another collapse, but rather imposes extra stability on the usual kernel.
	
	Thus multiplication has no separate kernel effect; instead, it forces
	\[
	\ker(\varphi)
	\]
	to be closed under multiplication by arbitrary elements of \(R\), which is exactly why the kernel is an ideal.
	
	\paragraph{Kernel and injectivity.}
	Suppose
	\[
	\varphi:A\to B
	\]
	is a homomorphism between algebraic structures. Then, in general, the kernel is defined by
	\[
	\ker(\varphi)=\{a\in A:\varphi(a)=e_B\},
	\]
	where \(e_B\) denotes the neutral element in the codomain with respect to the operation being preserved. Thus the symbol depends on the structure: in groups written multiplicatively the identity is \(1\) or \(e\); in groups written additively the identity is \(0\); in vector spaces or modules under addition the identity is \(0\); and in rings, when a ring homomorphism is viewed through its additive group structure, the relevant identity for the kernel is again \(0\). Hence the kernel is always about collapse to the identity, not about the particular symbol \(0\) in itself.
	
	For a linear map
	\[
	T:V\to W,
	\]
	the underlying operation in \(W\) is vector addition, whose neutral element is the zero vector. Therefore
	\[
	\ker(T)=\{v\in V:T(v)=0\}.
	\]
	Here \(0\) means the additive identity of \(W\), not an arbitrary number. That is why in linear algebra the kernel is naturally described as the set of vectors sent to zero. These are exactly the elements that become invisible under the map.
	
	If
	\[
	\varphi:G\to H
	\]
	is a group homomorphism, then
	\[
	\ker(\varphi)=\{g\in G:\varphi(g)=e_H\}.
	\]
	If \(H\) is written multiplicatively, this becomes
	\[
	\ker(\varphi)=\{g\in G:\varphi(g)=1\},
	\]
	whereas if \(H\) is written additively, it becomes
	\[
	\ker(\varphi)=\{g\in G:\varphi(g)=0\}.
	\]
	So the image of the kernel may be written as \(e\), \(1\), or \(0\), depending on how the codomain operation is expressed.
	
	The kernel controls injectivity because, for a homomorphism, the condition
	\[
	\varphi(a)=\varphi(b)
	\]
	means that the map cannot distinguish \(a\) from \(b\). In the group setting,
	\[
	\varphi(a)=\varphi(b)
	\iff \varphi(a)\varphi(b)^{-1}=e
	\iff \varphi(ab^{-1})=e
	\iff ab^{-1}\in\ker(\varphi).
	\]
	Thus two elements have the same image exactly when they differ by an element of the kernel. In this precise sense, the kernel measures the failure of injectivity. If the only element sent to the identity is the identity itself, then no nontrivial collapse occurs. Hence
	\[
	\varphi \text{ is injective } \iff \ker(\varphi)=\{e_A\}.
	\]
	
	For a linear map
	\[
	T:V\to W,
	\]
	one has
	\[
	T(u)=T(v)
	\iff T(u-v)=0
	\iff u-v\in\ker(T).
	\]
	Therefore, if
	\[
	\ker(T)=\{0\},
	\]
	then
	\[
	T(u)=T(v)\implies u-v=0\implies u=v,
	\]
	so \(T\) is injective. Conversely, if \(T\) is injective and \(T(v)=0\), then also
	\[
	T(v)=T(0),
	\]
	and injectivity implies
	\[
	v=0.
	\]
	Therefore
	\[
	\ker(T)=\{0\}.
	\]
	
	A clean formulation is the following: the kernel of a homomorphism is the set of domain elements sent to the identity element of the codomain; therefore it records exactly which elements are collapsed by the map, and hence it measures injectivity. So the more exact version of the intuition is not simply that the kernel consists of elements that collapse to \(0\), but rather that it consists of elements sent to the codomain identity, that is, those elements that become indistinguishable from the neutral element under the map. In linear algebra this means collapse to \(0\); in multiplicative groups it means collapse to \(e\) or \(1\); and in additive groups it means collapse to \(0\). The final structural point is therefore that the kernel is not specifically about \(0\), but about the neutral element of the codomain structure. In linear maps that neutral element is \(0\), while in group homomorphisms it may be \(e\) or \(1\). This is the general algebraic idea behind the notion of kernel.
	
	We need the kernel because, once a map \(T:V\to W\) is established, we want to understand what structural information the map preserves and what it destroys. The kernel records precisely the part of the domain that the map cannot detect. More explicitly, it serves several fundamental purposes. First, it determines whether the map is injective, since
	\[
	T(u)=T(v)\iff T(u-v)=0\iff u-v\in\ker(T).
	\]
	Thus, if
	\[
	\ker(T)=\{0\},
	\]
	then the only vector sent to zero is the zero vector, and therefore distinct inputs cannot collapse to the same output; in this case, \(T\) is injective. In this sense, the kernel measures the failure of injectivity. Second, it describes which directions in the domain are collapsed by the map: if \(k\in\ker(T)\), then for every \(u\in V\),
	\[
	T(u+k)=T(u),
	\]
	so moving along a kernel direction changes the input without changing the output. Hence the map identifies whole families of vectors as equivalent from the point of view of the codomain. Third, the kernel describes the effective content of the map: a map may be defined on a large space, but if a large subspace lies in the kernel, then much of the domain structure is lost under the mapping. Fourth, together with the image, it provides the main structural decomposition of the map: the kernel records what is lost, while the image records what is reached. In finite-dimensional linear algebra, this is expressed by the rank--nullity formula
	\[
	\dim(V)=\dim(\ker(T))+\dim(\operatorname{Im}(T)).
	\]
	Thus, the dimension of the domain splits into the dimension of what is annihilated and the dimension of what effectively passes through the map. Fifth, the kernel leads naturally to quotient spaces: since vectors differing by a kernel element have the same image, the map depends not on each vector individually, but on its class modulo the kernel. In this sense, the kernel indicates exactly what must be quotiented out in order to obtain the essential injective form of the map. Therefore, the deepest structural role of the kernel is to identify the ambiguity introduced by the map, measure its loss of information, test injectivity, and prepare the passage to quotient structures. In compact form: once a map between two vector spaces is established, the kernel tells us which inputs become indistinguishable under the map.
	
\end{document}